\documentclass[11pt,a4paper]{article}

\usepackage[utf8]{inputenc}
\usepackage[T1]{fontenc}
\usepackage[spanish,es-tabla]{babel}
\usepackage{lmodern}
\usepackage[margin=2.3cm]{geometry}
\usepackage{booktabs}
\usepackage{tabularx}
\usepackage{array}
\usepackage{microtype}
\usepackage{xcolor}
\usepackage[colorlinks=true,linkcolor=blue!50!black,urlcolor=blue!50!black]{hyperref}

\setlength{\parskip}{0.45em}
\setlength{\parindent}{0pt}
\newcolumntype{Y}{>{\raggedright\arraybackslash}X}

\title{\textbf{E01 · Documento ejecutivo}\\[0.3em]
\large Frecuencias inteligentes para SOBUSA: demanda prevista, despachos recomendados y alertas}
\author{Roger Daza \quad Matteo Del Sole \quad Rafael Romero\\[0.2em]\small Universidad del Norte · Reto Extreme Technologies -- SOBUSA}
\date{29 de octubre de 2026}

\begin{document}
\maketitle

\section*{El problema}

Hoy la frecuencia de despacho de cada ruta se fija con la programación vigente y la experiencia de los despachadores. Cuando la demanda cambia en una franja, el ajuste llega tarde: hay horas con buses llenos y pasajeros esperando, y otras con buses casi vacíos que consumen flota, conductores y combustible. La pregunta del reto es concreta: con los datos históricos y recientes, \emph{¿cada cuánto despachar buses para optimizar flota y recaudo, cuántos pasajeros vendrán y cuántos vehículos se necesitan?}

\section*{Qué entregamos}

Un prototipo en funcionamiento que, cada hora y sin intervenir la base de datos del operador:

\begin{enumerate}
  \setlength{\itemsep}{0.15em}
  \item \textbf{Lee los datos nuevos} de programación y de los sensores de conteo de pasajeros, en modo de solo lectura, y valida su calidad.
  \item \textbf{Pronostica las subidas} de las próximas dos horas por ruta, en franjas de 15 minutos, con un rango de confianza.
  \item \textbf{Recomienda la frecuencia}: buses por hora, intervalo entre despachos y vehículos necesarios, balanceando costo de operación, espera de los pasajeros, sobreocupación y estabilidad del servicio, dentro de los límites de flota y headway.
  \item \textbf{Genera alertas explicables}: diez tipos, desde ``la demanda superará la capacidad'' hasta ``despachos que no salieron'' o ``el modelo se degradó''. Cada alerta muestra el valor real, el esperado, el umbral y qué hacer.
  \item \textbf{Aprende de forma controlada}: vigila su propio error y reentrena solo cuando el modelo nuevo demuestra ser mejor que el vigente, con posibilidad de volver atrás.
\end{enumerate}

Todo se consulta en un tablero web y en una API documentada. Cada cifra queda trazada al dato, al modelo y a la regla que la produjo.

\section*{Resultados}

Se trabajó con tres rutas (SILENCIO, CRA 54 y CL 72) y 176 millones de lecturas de sensores de 2023 a 2025.

\begin{table}[h]
\centering\small
\begin{tabularx}{\textwidth}{lYY}
\toprule
 & \textbf{Referencia sin modelo} & \textbf{Prototipo} \\
\midrule
Error del pronóstico (WAPE), prueba mayo--noviembre 2025 & 14,9\,\% (promedio de las 4 semanas previas) & \textbf{10,1\,\%, 10,5\,\% y 10,8\,\%} a 30, 60 y 120 min \\
Error en días especiales (festivos, carnaval, Semana Santa) & 99\,\% & \textbf{14--16\,\%} \\
Rango de confianza del 80\,\% & --- & cubre \textbf{81,6--82,5\,\%} de los casos \\
Consultas del tablero (p95, 10 usuarios simultáneos) & límite pedido: 2 s & \textbf{146 ms} \\
Datos nuevos procesados & mínimo pedido: 100.000 registros & \textbf{808.345 registros en 149 s}, sin duplicados \\
\bottomrule
\end{tabularx}
\end{table}

\textbf{La recomendación propone redistribuir la oferta, no recortarla.} Se recalcularon las 748 horas-ruta del 1 al 14 de diciembre de 2025:

\begin{itemize}
  \setlength{\itemsep}{0.1em}
  \item \textbf{Total de buses.} El modelo propone en promedio 8,4 buses por hora frente a 8,5 operados.
  \item \textbf{SILENCIO.} Pasaría de 7,7 a 6,6 buses por hora: hoy opera con holgura.
  \item \textbf{CL 72.} Pasaría de 4,6 a 5,5 buses por hora, porque hoy opera por debajo del intervalo máximo de 10 minutos configurado como nivel de servicio.
  \item \textbf{CRA 54.} Queda prácticamente igual (12,8 frente a 12,7).
\end{itemize}

La ocupación esperada con la recomendación tiene mediana de 65\,\%, el piso del rango objetivo de 65--85\,\%.

\section*{Valor para SOBUSA}

\begin{itemize}
  \setlength{\itemsep}{0.15em}
  \item \textbf{Anticipación.} Saber con hasta dos horas de antelación cuántos pasajeros vienen, con un error cercano al 10\,\%, permite despachar antes de que se forme la fila.
  \item \textbf{Mejor uso de la flota.} Mover buses de donde sobran a donde faltan con la misma flota, y justificar cada cambio con números.
  \item \textbf{Control operativo.} Las alertas de despacho detectan horas en que una cuarta parte o más de los despachos programados no salió a tiempo. Con los datos de diciembre son unas 4 por día entre las tres rutas.
  \item \textbf{Gobierno.} Umbrales, pesos y modelos versionados en la base de datos, cambiables sin programar, con bitácora de quién cambió qué. La base del operador no se toca: solo se lee.
\end{itemize}

\section*{Limitaciones}

\begin{itemize}
  \setlength{\itemsep}{0.15em}
  \item \textbf{Capacidad y nivel de servicio por validar.} Usamos 50 pasajeros por bus (mediana observada en 2025), ocupación objetivo del 80\,\% e intervalo máximo de 10 minutos de día. En 7--20\,\% de los viajes la carga medida supera los 50 pasajeros, lo que sugiere que la capacidad real es mayor. Estos valores explican buena parte de las alertas de sobreocupación; la calibración de umbrales solo las redujo 14\,\%.
  \item \textbf{Pesos de la función objetivo.} El valor relativo de una hora de espera frente a una hora de bus son referencias del equipo. Definen cuánto servicio se cambia por costo, y deben fijarlos SOBUSA y Extreme Technologies.
  \item \textbf{Fuente de datos.} Los consolidados \texttt{timbradas\_*} entregados cubren solo ocho semanas de 2023 y tienen horas desplazadas. Por eso la demanda se reconstruyó desde los sensores. Los scripts de segmentación no se pudieron ejecutar con acceso de solo lectura; se replicó su lógica, con 99,7\,\% de coincidencia.
  \item \textbf{Notificaciones.} Correo y WhatsApp están implementados, pero sin credenciales del operador los envíos quedan simulados.
  \item \textbf{Alcance.} Tres rutas y datos hasta diciembre de 2025. No se probó en producción ni con el conjunto reservado del operador.
\end{itemize}

\section*{Recomendación}

\begin{enumerate}
  \setlength{\itemsep}{0.15em}
  \item \textbf{Piloto de 4 semanas en CL 72 y SILENCIO}, las rutas con mayor diferencia entre lo operado y lo recomendado. La recomendación es consultiva: el despachador decide. Se mide la espera, la ocupación y los despachos fallidos frente a las mismas semanas del año anterior.
  \item \textbf{Validar con operación}, antes del piloto, la capacidad por tipología de bus, el intervalo máximo por ruta y franja, y los pesos de la función objetivo. Son parámetros de la base de datos: cambiarlos no requiere desarrollo.
  \item \textbf{Entregar al prototipo} credenciales de correo y WhatsApp, la flota disponible por ruta y acceso a los consolidados completos, si existen. Así las alertas llegan a quien debe actuar y la recomendación respeta la flota real.
  \item \textbf{Extender a más rutas} cuando el piloto confirme el beneficio. La solución no depende de la ruta: basta con cargar su histórico y reentrenar.
\end{enumerate}

\vfill
\small\textit{Detalle técnico: E02 · Especificación técnica; integración con PostgreSQL: E07; código, pruebas y manuales en el repositorio (README).}

\end{document}
